\documentclass[10pt,a4paper]{amsart}
\usepackage{amsmath,amssymb,amsthm,graphicx}
\usepackage{verbatim}
\newenvironment{reference}{\comment}{\endcomment}
%\newenvironment{reference}{}{}
\newenvironment{slogan}{\comment}{\endcomment}
\newenvironment{history}{\comment}{\endcomment}

% For commutative diagrams we use Xy-pic
\usepackage[all]{xy}

% We use 2cell for 2-commutative diagrams.
\xyoption{2cell}
% Theorem environments.
%
\theoremstyle{plain}
\newtheorem{theorem}[subsection]{Theorem}
\newtheorem{proposition}[subsection]{Proposition}
\newtheorem{lemma}[subsection]{Lemma}

\theoremstyle{definition}
\newtheorem{definition}[subsection]{Definition}
\newtheorem{example}[subsection]{Example}
\newtheorem{exercise}[subsection]{Exercise}
\newtheorem{situation}[subsection]{Situation}

\theoremstyle{remark}
\newtheorem{remark}[subsection]{Remark}
\newtheorem{remarks}[subsection]{Remarks}

\numberwithin{equation}{subsection}

% Macros
%
\def\lim{\mathop{\mathrm{lim}}\nolimits}
\def\colim{\mathop{\mathrm{colim}}\nolimits}
\def\Spec{\mathop{\mathrm{Spec}}}
\def\Hom{\mathop{\mathrm{Hom}}\nolimits}
\def\Ext{\mathop{\mathrm{Ext}}\nolimits}
\def\SheafHom{\mathop{\mathcal{H}\!\mathit{om}}\nolimits}
\def\SheafExt{\mathop{\mathcal{E}\!\mathit{xt}}\nolimits}
\def\Sch{\mathit{Sch}}
\def\Mor{\mathop{\mathrm{Mor}}\nolimits}
\def\Ob{\mathop{\mathrm{Ob}}\nolimits}
\def\Sh{\mathop{\mathit{Sh}}\nolimits}
\def\NL{\mathop{N\!L}\nolimits}
\def\CH{\mathop{\mathrm{CH}}\nolimits}
\def\proetale{{pro\text{-}\acute{e}tale}}
\def\etale{{\acute{e}tale}}
\def\QCoh{\mathit{QCoh}}
\def\Ker{\mathop{\mathrm{Ker}}}
\def\Im{\mathop{\mathrm{Im}}}
\def\Coker{\mathop{\mathrm{Coker}}}
\def\Coim{\mathop{\mathrm{Coim}}}

% Boxtimes
%
\DeclareMathSymbol{\boxtimes}{\mathbin}{AMSa}{"02}

%
% Macros for moduli stacks/spaces
%
\def\QCohstack{\mathcal{QC}\!\mathit{oh}}
\def\Cohstack{\mathcal{C}\!\mathit{oh}}
\def\Spacesstack{\mathcal{S}\!\mathit{paces}}
\def\Quotfunctor{\mathrm{Quot}}
\def\Hilbfunctor{\mathrm{Hilb}}
\def\Curvesstack{\mathcal{C}\!\mathit{urves}}
\def\Polarizedstack{\mathcal{P}\!\mathit{olarized}}
\def\Complexesstack{\mathcal{C}\!\mathit{omplexes}}
% \Pic is the operator that assigns to X its picard group, usage \Pic(X)
% \Picardstack_{X/B} denotes the Picard stack of X over B
% \Picardfunctor_{X/B} denotes the Picard functor of X over B
\def\Pic{\mathop{\mathrm{Pic}}\nolimits}
\def\Picardstack{\mathcal{P}\!\mathit{ic}}
\def\Picardfunctor{\mathrm{Pic}}
\def\Deformationcategory{\mathcal{D}\!\mathit{ef}}
\usepackage{cleveref}
\crefname{originalreference}{original reference}{original references}
\Crefname{originalreference}{Original reference}{Original references}
\makeatletter
\newlabel{lemma-algebraic-goes-up}{{09GF}{source}{}{sourcepointer-1}{}}
\newlabel{lemma-algebraic-goes-up@cref}{{[originalreference][0][]09GF}{[1][0][]source}{}{}{}}
\makeatother
\title{Every algebraic extension maps into any algebraic closure of the base field}
\author{The Stacks Project Authors}
\date{stacks-project:fields.tex:local-snapshot-20261012; GNU FDL 1.2 or later}
\begin{document}
\maketitle
\noindent\textit{Editorial note.} All delivered mathematical text is the author's own, copied from the pinned Stacks Project chapter file fields.tex; no statement or proof step was rewritten, reordered or removed.\par\medskip
\noindent\textit{Editorial note.} The counted proof is the author's complete argument: The source line 1257 contains the author's typo 'The alpha is algebraic' for 'Then alpha is algebraic'; preserved verbatim per the no-content-edits rule.\par\medskip
\noindent\textit{Editorial note.} The label lemma-map-into-algebraic-closure names the counted statement, and the proof environment immediately following it in the source is delivered verbatim as its proof.\par\medskip
\noindent\textit{Editorial note.} The counted statement is carried by the author's own lemma environment in the chapter source.\par\medskip
\noindent\textit{Allowed standard prerequisites (editorial).}\begin{itemize}
\item Context span 1169-1177: definition of algebraic closure
\item Same chapter: lemma-algebraic-goes-up (elements of an algebraic tower are algebraic down the tower, source cited in the proof); Zorn's lemma
\end{itemize}\medskip
\ifdefined\hypertarget\hypertarget{sourcepointer-1}{}\fi
\noindent\textit{Source reference.} Original reference 09GF: sources/stacks-fields/pinned-original-source.tex, source lines 860--860. Prerequisite pointer invoked at 1258; it names a standard fact of the same chapter that the counted argument uses, and every step of the counted proof built on it is retained verbatim in the delivered body.\par\medskip
\begin{definition}
\label{definition-algebraic-closure}
Let $F$ be a field. An {\it algebraic closure} of $F$ is a field
$\overline{F}$ containing $F$ such that:
\begin{enumerate}
\item $\overline{F}$ is algebraic over $F$.
\item $\overline{F}$ is algebraically closed.
\end{enumerate}
\end{definition}

\begin{lemma}
\label{lemma-map-into-algebraic-closure}
Let $F$ be a field. Let $\overline{F}$ be an algebraic closure of $F$.
Let $M/F$ be an algebraic extension. Then there is a morphism of
$F$-extensions $M \to \overline{F}$.
\end{lemma}

\begin{proof}
Consider the set $I$ of pairs $(E, \varphi)$ where $F \subset E \subset M$
is a subextension and $\varphi : E \to \overline{F}$ is a morphism of
$F$-extensions. We partially order the set $I$ by declaring
$(E, \varphi) \leq (E', \varphi')$ if and only if $E \subset E'$ and
$\varphi'|_E = \varphi$. If $T = \{(E_t,  \varphi_t)\} \subset I$
is a totally ordered subset, then
$\bigcup \varphi_t : \bigcup E_t \to \overline{F}$ is an element of $I$.
Thus every totally ordered subset of $I$ has an upper bound.
By Zorn's lemma there exists a maximal element $(E, \varphi)$ in $I$.
We claim that $E = M$, which will finish the proof. If not, then
pick $\alpha \in M$, $\alpha \not \in E$. The $\alpha$ is algebraic
over $E$, see Lemma \ref{lemma-algebraic-goes-up}.
Let $P$ be the minimal polynomial of $\alpha$ over $E$.
Let $P^\varphi$ be the image of $P$ by $\varphi$ in $\overline{F}[x]$.
Since $\overline{F}$ is algebraically closed there is a root $\beta$
of $P^\varphi$ in $\overline{F}$. Then we can extend $\varphi$ to
$\varphi' : E(\alpha) = E[x]/(P) \to \overline{F}$ by mapping
$x$ to $\beta$. This contradicts the maximality of $(E, \varphi)$
as desired.
\end{proof}

\end{document}
